\chapter{Linear ODEs and the Fundamental Matrix}
\label{ch:differential-equations}

A fundamental matrix is a reusable computation depending only on the
coefficient matrix.  It propagates initial states, accumulates ordinary
forcing, and propagates delta inputs.  We construct it from iterated finite
integrals and factorial error bounds.

\section{Coefficient data and finite norms}
Fix a rational interval $[a,b]$.  The entries of $A(t)$ are supplied
rational-input computations with interval-continuity certificates and a
rational bound $\|A(t)\|\le M$.  Use the maximum norm for vectors and the
maximum absolute row-sum norm for matrices; finite algebra proves
$\|AB\|\le\|A\|\|B\|$ and $\|Av\|\le\|A\|\|v\|$.

Polynomial matrices are the first examples.  Separated rational functions
and the elementary computations of the preceding chapters provide others.
A finite set of coefficient breakpoints is handled by solving on each
piece and composing the resulting evolutions.

\section{Peano--Baker prefixes}
For rational $s,t\in[a,b]$, define
\[
 P_0(t,s)=I,\qquad P_{n+1}(t,s)=\int_s^t A(u)P_n(u,s)\,du,
 \qquad \Phi_N(t,s)=\sum_{n=0}^NP_n(t,s).
\]
For fixed $n$ these are finite nested integral computations.  Their
continuity data is obtained recursively from products and the
bounded-integrand endpoint estimate.  If $A$ is a rational-polynomial
matrix, all $P_n$ are computed by finite polynomial integration.
Matrix multiplication is kept in the displayed order; matrices at distinct
times are not assumed to commute.

For $h=|t-s|$, induction using polynomial integration gives
\[
 \|P_n(t,s)\|\le\frac{(Mh)^n}{n!}.
\]
The argument works with oriented integrals in either time direction.
If $z=M(b-a)$ and $N+2\ge2z$, every finite tail after $N$ is bounded by
\[
 \boxed{2\frac{z^{N+1}}{(N+1)!}.}
\]
Compute the finite prefix to an additional rational error budget and
stabilize the resulting boxes.  This constructs the matrix computation
$\Phi(t,s)$ uniformly on the interval.  A rational upper bound $B$ for
$E(M(b-a))$ bounds every $\|\Phi(t,s)\|$.

\begin{figure}[htbp]\centering
\peanoBakerWordsAnimation
\caption{Ordered matrix products and their finite integral prefixes.}
\end{figure}

The finite recursion, with its uniform tails, gives
\[
 \Phi(t,s)=I+\int_s^t A(u)\Phi(u,s)\,du.
\]
The accumulation estimate therefore constructs its derivative:
\[
 \partial_t\Phi(t,s)=A(t)\Phi(t,s),\qquad\Phi(s,s)=I.
\]
No general differentiation-of-a-limit assertion has been used.

\section{Uniqueness belongs to the integral construction}
Suppose a bounded computed vector $w$ satisfies
\[
 w(t)=\int_s^t A(u)w(u)\,du.
\]
Starting with any rational bound $\|w\|\le C$, repeated substitution gives
\[
 \|w(t)\|\le C\frac{(M|t-s|)^n}{n!}
\]
for every integer $n$.  The right side tends to zero with an explicit
factorial bound, so $w$ is equivalent to zero.

For the solutions considered here, the integral identity is proved by
construction.  We do not infer it from a bare equation at every rational
point; a function hidden across a cut could satisfy such a pointwise
equation and fail its integral identity.

The same uniqueness argument gives, for all rational $s,u,t$,
\[
 \boxed{\Phi(t,u)\Phi(u,s)=\Phi(t,s).}
\]
Both sides have the same integral equation and initial value when regarded
as functions of $t$.  The identity follows for every starting location by
subtracting their integral equations there.  In particular
$\Phi(s,t)$ is the inverse of $\Phi(t,s)$; no determinant inversion is
needed to construct it.

The endpoint estimate gives
$\|\Phi(u,s)-I\|\le MB|u-s|$.  Together with composition it yields
\[
 \|\Phi(t,s)-\Phi(t,u)\|\le MB^2|s-u|.
\]
This supplies continuity in the source time as well as the target time,
and permits computed times with domain enclosures.

\section{Ordinary forcing and variation of constants}
For a supplied interval-continuous forcing $r$, put
\[
 y(t)=\Phi(t,a)y_a+\int_a^t\Phi(t,s)r(s)\,ds.
\]
All entries of this integrand have the continuity and boundedness data just
proved.  The finite integral recursions show
\[
 y(t)=y_a+\int_a^t\bigl(A(u)y(u)+r(u)\bigr)\,du.
\]
One verifies the double-integral rearrangement on finite subdivisions of
the triangle $a\le s\le u\le t$.  The cells meeting its diagonal have
total area bounded by a constant times $(b-a)\mesh(P)$; the integrands
have rational bounds, and the other cells are finite iterated rectangles.
This proves the needed interchange without a general measure-theoretic
Fubini theorem.

Accumulation gives $Dy=Ay+r$, and the preceding uniqueness estimate proves
uniqueness among the certified integral solutions.  Thus the displayed
formula is the solution computation, not a formula justified by an
unconstructed classical solution.

The same estimate controls numerical replacements.  If a proposed
computation $\widetilde y$ has a certified integral identity with initial
error at most $e$ and additional forcing (residual) bounded by $\rho$,
then variation of constants gives
\[
 \|\widetilde y(t)-y(t)\|\le B\bigl(e+(b-a)\rho\bigr).
\]
A merely sampled residual would not suffice: its bound must hold on the
cells, with the integral identity certified for that computation.

\section{Delta inputs are finite state updates}
For a finite list of impulse times $a<\tau_j<b$ and vectors $c_j$, solve
\[
 Dy=A(t)y+r(t)+\sum_jc_j\delta_{\tau_j}.
\]
Initially take the times rational.  For times away from the impulses the
solution is
\[
 \boxed{y(t)=\Phi(t,a)y_a+\int_a^t\Phi(t,s)r(s)\,ds
                  +\sum_{a<\tau_j<t}\Phi(t,\tau_j)c_j.}
\]
Its one-sided states satisfy
$y(\tau_j+)=y(\tau_j-)+c_j$ because $\Phi(\tau_j,\tau_j)=I$.
The piecewise integration-by-parts rule of Chapter~\ref{ch:impulses}
proves exactly the displayed equation on tests.  The ordinary branches
have the needed integration identities from their integral equations and
the corresponding finite product rule with a smooth test.

One computation evaluates the formula directly.  Another evolves to each
impulse, adds its vector, and continues.  The composition law proves their
equivalence.  A computed nonrational impulse location can also be used
when ordering and separated evaluation times are supplied; one-sided
states are retained instead of forcing an undecidable comparison at the
location itself.

Define the causal matrix kernel
\[
 G(t,s)=\begin{cases}0,&t<s,\\\Phi(t,s),&t>s.\end{cases}
\]
Its jump is $I$, and hence
\[
 \boxed{(D_t-A(t))G(t,s)=I\delta_s.}
\]
Each column is the response to a unit impulse in one state coordinate.
This identity is the reason for including delta inputs with the fundamental
matrix.

For a rectangular pulse of width $\varepsilon$, after the pulse ends the
response is $\varepsilon^{-1}\int_s^{s+\varepsilon}\Phi(t,u)c\,du$.
The source-time bound gives its difference from $\Phi(t,s)c$ as at most
$MB^2\|c\|\varepsilon/2$.  This is agreement between the state-update
construction and narrowing ordinary forcing, not uniform convergence
through a discontinuity.

\section{Three explicit kernels}
\subsection{The free second derivative}
For the state $(y,y')$ and matrix
$A=\left(\begin{smallmatrix}0&1\\0&0\end{smallmatrix}\right)$,
\[
 \Phi(t,s)=\begin{pmatrix}1&t-s\\0&1\end{pmatrix}.
\]
An impulse $(0,1)^{\mathsf T}$ gives $y=(t-s)_+$, so $D^2y=\delta_s$.
The function $|t-s|$ solves the equation with $2\delta_s$ but differs
from its zero-past solution by an affine function.  Initial data decides
that homogeneous ambiguity.

\subsection{The oscillator}
For $y''+\pi^2y=0$,
\[
 \Phi(t,s)=
 \begin{pmatrix}
 C(t-s)&S(t-s)/\pi\\
 -\pi S(t-s)&C(t-s)
 \end{pmatrix}.
\]
Its entries, derivatives, and composition are already certified by the
circle and FTC chapters.  The impulse in the velocity coordinate gives
\[
 g_s(t)=\begin{cases}0,&t<s,\\S(t-s)/\pi,&t>s,\end{cases}
 \qquad (D^2+\pi^2)g_s=\delta_s.
\]
The position is continuous and the velocity jumps by $1$.  A position jump
would instead introduce a derivative of delta into the second derivative.

\subsection{A variable coefficient: the Airy equation}
For $y''-ty=0$, substitution of $y=\sum a_nt^n$ gives
\[
 a_2=0,\qquad a_{n+3}=\frac{a_n}{(n+3)(n+2)}.
\]
Take $u(0)=1,u'(0)=0$ and $v(0)=0,v'(0)=1$.  Thus
\[
 u(t)=1+t^3/6+t^6/180+\cdots,\qquad
 v(t)=t+t^4/12+t^7/504+\cdots.
\]
On $|t|\le R$, each of the three coefficient strands eventually has
successive absolute-term ratio at most $1/2$, by taking
$(n+3)(n+2)\ge2R^3$.  This gives geometric tail bounds; differentiated
strands have the same property after increasing the index.  Thus $u,v$
and their derivatives are actual series computations on every bounded
interval.

Coefficient multiplication proves $uv'-u'v=1$: its formal derivative is
$uv''-u''v=0$, so all nonconstant coefficients vanish, and its constant
coefficient is $1$.  Tails justify the identity for the computations.
Consequently
\[
 g_s(t)=\begin{cases}
 0,&t<s,\\
 u(s)v(t)-v(s)u(t),&t>s
 \end{cases}
\]
is continuous at $s$ and its first derivative jumps by $1$.  Hence
\[
 (D^2-t)g_s=\delta_s.
\]
This constructs a variable-coefficient impulse response by rational
recurrences and finite determinant algebra.  It is equivalent to the
Peano--Baker computation by integral-equation uniqueness.

\section{The FTC as the zero-coefficient case}
When $A=0$, the fundamental matrix is $I$.  Ordinary forcing and the
finite atoms give
\[
 y(t)-y(a)=\int_a^t r(s)\,ds+\sum_{a<\tau_j<t}c_j,
\]
with endpoints away from impulses.  Thus integration of a derivative is
inversion of $D$, up to a constant; linear ODE evolution is inversion of
$D-A$, up to a homogeneous solution.  Both are identities between explicit
computations with initial data, finite updates, and shrinking errors.
